\section{Поля разложения}

Пусть $f\in F[X]$~--- ненулевой многочлен с коэффициентами в $F$. Говорят, что поле $E\supset F$ \emph{расщепляет} $f$ (или что
$f$ \emph{раскладывается} в $E$), если $f$ раскладывается в $E[X]$ на линейные множители:
\[
f=a\prod_{i=1}^m(X-\alpha_i),\qquad a\in F,\ \alpha_i\in E .
\]
Если, кроме того, $E$ порождено над $F$ корнями $f$, т.\,е.\ $E=F[\alpha_1,\dots,\alpha_m]$, то $E$ называется \emph{полем
разложения} многочлена $f$. (Кольцо $F[\alpha_1,\dots,\alpha_m]$, порождённое алгебраическими элементами, само является
полем и конечно над $F$~--- глава~1, замечание после \rf{1.30}.)

\begin{remark}[2.4e]
(1) Многочлены $\prod f_i^{m_i}$ ($m_i\ge1$) и $\prod f_i$ имеют одни и те же поля разложения: у них одни и те же корни, и
каждый из них раскладывается на линейные множители там же, где и другой (если раскладывается $\prod f_i^{m_i}$, то по
однозначности разложения раскладывается каждый $f_i$).

(2) Чтобы проверить, что $E$ расщепляет $f$, $\deg f=n$, достаточно найти в $E$ корни $\alpha_1,\dots,\alpha_{n-1}$ с
учётом кратности, т.\,е.\ разложение $f=(X-\alpha_1)\cdots(X-\alpha_{n-1})\,g$ в $E[X]$: тогда $\deg g=1$, и последний
корень тоже в $E$. Его можно указать явно: если $f=aX^n+a_1X^{n-1}+\dots$, то сумма корней $\sum\alpha_i=-a_1/a$ лежит в $F$.

(3) Если $E$ расщепляет $f$ и $f=a\prod(X-\alpha_i)$ в $E[X]$, то подполе $F[\alpha_1,\dots,\alpha_m]\subset E$~--- поле
разложения $f$, и это единственное поле разложения $f$, содержащееся в $E$: всякое такое поле содержит все корни $f$ в
$E$ (по однозначности разложения в $E[X]$ корни $f$ в $E$~--- это в точности $\alpha_i$) и порождено ими.
\end{remark}

\begin{example}[2.5]
Пусть $f=aX^2+bX+c\in\Q[X]$, $a\ne0$, и $\alpha=\sqrt{b^2-4ac}\in\CC$. Подполе $\Q[\alpha]$ поля $\CC$~--- поле разложения $f$.
Действительно, в $\Q[\alpha][X]$
\[
f=a\Bigl(X-\frac{-b-\alpha}{2a}\Bigr)\Bigl(X-\frac{-b+\alpha}{2a}\Bigr),
\]
так что корни $x_1=\frac{-b-\alpha}{2a}$, $x_2=\frac{-b+\alpha}{2a}$ лежат в $\Q[\alpha]$; обратно, $\alpha=a(x_2-x_1)\in\Q[x_1,x_2]$.
Значит, $\Q[\alpha]=\Q[x_1,x_2]$. Заметим, что $\alpha^2=a^2(x_2-x_1)^2=b^2-4ac$~--- \emph{дискриминант} $f$, а $\alpha$~---
квадратный корень из него (ср.\ \rf{2.6}).
\end{example}

\begin{example}[2.6]
Пусть $f=X^3+aX^2+bX+c\in\Q[X]$ неприводим и $\alpha_1,\alpha_2,\alpha_3$~--- его корни в $\CC$. Тогда
$\Q[\alpha_1,\alpha_2,\alpha_3]=\Q[\alpha_1,\alpha_2]$ (так как $\alpha_3=-a-\alpha_1-\alpha_2$)~--- поле разложения $f$. Далее,
$[\Q[\alpha_1]:\Q]=3$, а $[\Q[\alpha_1,\alpha_2]:\Q[\alpha_1]]$ равно $1$ или $2$: в $\Q[\alpha_1][X]$ имеем
$f=(X-\alpha_1)g$ с квадратичным $g$, и $\alpha_2$~--- корень $g$; степень равна $1$, если $g$ раскладывается над
$\Q[\alpha_1]$, и $2$, если $g$ неприводим. Итак, $[\Q[\alpha_1,\alpha_2]:\Q]=3$ или $6$. Позже (\rf{4.7}) мы увидим, что
степень равна $3$ тогда и только тогда, когда дискриминант $f$ является квадратом в $\Q$; например, дискриминант
$X^3+bX+c$ равен $-4b^3-27c^2$.

\emph{Степень $6$:} $f=X^3+10X+1$ (рациональных корней нет: $f(\pm1)=12,-10$), дискриминант $-4027$~--- не квадрат.
Здесь можно обойтись и без \rf{4.7}: так как $f'=3X^2+10>0$, у $f$ ровно один вещественный корень $\alpha_1$, и
$\Q[\alpha_1]\subset\R$ не содержит невещественных корней $\alpha_2,\alpha_3$; поэтому $g$ неприводим над $\Q[\alpha_1]$, и
степень равна $6$.

\emph{Степень $3$:} $f=X^3-3X+1$ (рациональных корней нет: $f(\pm1)=-1,3$; дискриминант $81=9^2$). Если $\alpha$~--- корень, то $\alpha^2-2$ и $-\alpha^2-\alpha+2$~---
тоже корни (это проверяется делением с остатком на $f$), и все три различны (их записи в базисе $1,\alpha,\alpha^2$
различны). Итак, $\Q[\alpha]$~--- поле разложения, степень $3$. (В $\CC$: $\alpha=2\cos\frac{2\pi}9$, остальные корни
$2\cos\frac{4\pi}9$, $2\cos\frac{8\pi}9$.)
\end{example}

\begin{proposition}[2.7]
Каждый ненулевой многочлен $f\in F[X]$ имеет поле разложения $E_f$, и
\[
[E_f:F]\le(\deg f)!\,.
\]
\end{proposition}

\begin{proof}
Индукция по $n=\deg f$, одновременно для всех полей $F$. Если $n\le1$, то $f$ уже раскладывается над $F$, и $E_f=F$.
Пусть $n\ge2$. Возьмём унитарный неприводимый делитель $g_1$ многочлена $f$ в $F[X]$ и его поле корня
$F_1=F[\alpha_1]=F[X]/(g_1)$ (\rf{1.25}); тогда $f(\alpha_1)=0$ и $[F_1:F]=\deg g_1\le n$. В $F_1[X]$ имеем
$f=(X-\alpha_1)f_1$, $\deg f_1=n-1$. По предположению индукции (для поля $F_1$) у $f_1$ есть поле разложения
$E=F_1[\alpha_2,\dots,\alpha_n]$ с $[E:F_1]\le(n-1)!$. Тогда $f=(X-\alpha_1)f_1$ раскладывается в $E[X]$ на линейные
множители, и
\[
E=F_1[\alpha_2,\dots,\alpha_n]=F[\alpha_1][\alpha_2,\dots,\alpha_n]=F[\alpha_1,\alpha_2,\dots,\alpha_n]
\]
порождено над $F$ корнями $f$. Значит, $E$~--- поле разложения $f$ над $F$, и по \rf{1.20}
$[E:F]=[E:F_1][F_1:F]\le(n-1)!\cdot n=n!$.
\end{proof}

\par\noindent\begin{minipage}{\linewidth}
\begin{center}
\begin{tikzpicture}[x=1cm,y=1cm]
\node (F) at (0,0) {$F$};
\node at (0.9,0) {$\subset$};
\node (F1) at (2.2,0) {$F[\alpha_1]$};
\node at (3.5,0) {$\subset$};
\node (F2) at (5.0,0) {$F[\alpha_1,\alpha_2]$};
\node at (6.6,0) {$\subset\ \cdots\ \subset$};
\node (Fn) at (9.3,0) {$F[\alpha_1,\dots,\alpha_n]=E_f$};
\draw[decorate,decoration={brace,mirror,amplitude=5pt},cold] (-0.2,-0.3) -- node[below=6pt,font=\footnotesize]{$\le n$} (2.9,-0.3);
\draw[decorate,decoration={brace,mirror,amplitude=5pt},cold] (3.0,-0.3) -- node[below=6pt,font=\footnotesize]{$\le n-1$} (6.0,-0.3);
\draw[decorate,decoration={brace,mirror,amplitude=5pt},cold] (7.2,-0.3) -- node[below=6pt,font=\footnotesize]{$\le 1$} (11.2,-0.3);
\draw[decorate,decoration={brace,amplitude=5pt},cnew] (-0.2,0.3) -- node[above=6pt,font=\footnotesize]{$[E_f:F]\le n\,(n-1)\cdots1=n!$} (11.2,0.3);
\end{tikzpicture}
\end{center}
\figcap{Построение поля разложения: на каждом шаге присоединяется корень неприводимого делителя оставшегося
множителя; степень $i$-го шага не больше $n-i+1$.}
\end{minipage}\par\medskip

Обратите внимание, почему построенное поле \emph{порождено} корнями $f$: на каждом шаге поле корня $F_{i-1}[X]/(g_i)$
порождено над $F_{i-1}$ одним элементом $\alpha_i$~--- классом $X$ (всякое поле, содержащее $F_{i-1}$ и корень $\gamma$
многочлена $g_i$, содержит подполе $F_{i-1}[\gamma]$, изоморфное этому полю корня, \rf{2.1}(b)), поэтому
по индукции $F_i=F[\alpha_1,\dots,\alpha_i]$. Существование алгебраического замыкания при этом не используется: нужен лишь
тот факт, что у непостоянного многочлена есть корень в некотором конечном расширении (глава~1, после \rf{1.25}).

\begin{remark}[2.7a]
Поле разложения зависит от основного поля. Пусть $F\subset L$, $f\in F[X]$, и $\Om_L$~--- поле разложения $f$,
рассматриваемого как многочлен из $L[X]$. Тогда $\Om_L$ расщепляет $f$, и по \rf{2.4e}(3) подполе $F[\alpha_1,\dots,\alpha_m]$,
порождённое над $F$ корнями $f$ в $\Om_L$,~--- поле разложения $f$ над $F$. Но само $\Om_L$ над $F$ полем разложения, как
правило, не является и может быть даже бесконечным над $F$. Например, для $f=X^3-2$, $F=\Q$, $L=\R$: $\Om_\R=\CC=\R[\omega]$,
а поле разложения над $\Q$~--- $\Q[\sqrt[3]2,\omega]$ степени $6$ (\rf{2.4c}(2)).
\end{remark}

\begin{example}[2.8]
Пусть $p$ простое и $f=(X^p-1)/(X-1)=X^{p-1}+\dots+X+1\in\Q[X]$. Если $\zeta$~--- один корень $f$ (в каком-либо
расширении), то $\zeta^p=1$ и $\zeta\ne1$ (так как $f(1)=p\ne0$), т.\,е.\ $\zeta$ имеет порядок $p$ в мультипликативной
группе. Поэтому $\zeta,\zeta^2,\dots,\zeta^{p-1}$~--- $p-1$ различных корней $X^p-1$, отличных от $1$, т.\,е.\ корней $f$, а
больше $p-1=\deg f$ корней у $f$ нет. Итак, остальные корни~--- это $\zeta^2,\dots,\zeta^{p-1}$, и поле разложения $f$~---
$\Q[\zeta]$ (в $\CC$: $\zeta=e^{2\pi i/p}$). Так как $f$ неприводим (\rf{1.16a}(c)), $[\Q[\zeta]:\Q]=p-1$.
\end{example}

\begin{example}[2.9]
Пусть $\operatorname{char}F=p\ne0$ и $f=X^p-X-a\in F[X]$. Если $\alpha$~--- корень $f$ в некотором расширении $F$, то для
$k\in\F_p\subset F$
\[
(\alpha+k)^p-(\alpha+k)-a=\alpha^p+k^p-\alpha-k-a=(\alpha^p-\alpha-a)+(k^p-k)=0,
\]
так как $(x+y)^p=x^p+y^p$ и $k^p=k$ (малая теорема Ферма). Элементы $\alpha,\alpha+1,\dots,\alpha+p-1$ попарно различны
($\alpha+i=\alpha+j\Rightarrow i=j$ в $\F_p$), и их $p=\deg f$ штук, поэтому
\[
f=\prod_{k\in\F_p}(X-\alpha-k),
\]
и поле разложения $f$~--- это $F[\alpha]$. Из этого описания корней можно вывести, что $f$ либо раскладывается над $F$
на линейные множители, либо неприводим (многочлен Артина--Шрайера; \rf{2.22b}, ср.\ \rf{5.32c});
например, над $\F_5$ многочлен $X^5-X$ раскладывается, а $X^5-X-1$ неприводим.
\end{example}

\begin{example}[2.10]
Пусть $f=X^n-a\in F[X]$, $a\ne0$, и $\alpha$~--- корень $f$ в некотором расширении $E$ поля $F$, расщепляющем $X^n-1$:
\[
X^n-1=\prod_{i=1}^n(X-\zeta_i)\qquad\text{в }E[X]
\]
(здесь $\zeta_i$~--- <<корни $n$-й степени из $1$>>, перечисленные с кратностями; такое $E$ есть по \rf{2.7},
применённому к полю $F[\alpha]$). Тогда
\[
X^n-a=\alpha^n\Bigl(\bigl(\tfrac X\alpha\bigr)^n-1\Bigr)=\alpha^n\prod_{i=1}^n\Bigl(\tfrac X\alpha-\zeta_i\Bigr)=\prod_{i=1}^n(X-\zeta_i\alpha),
\]
т.\,е.\ \emph{корни $X^n-a$~--- это $\zeta_i\alpha$}, где $\zeta_i$ пробегает корни $n$-й степени из $1$. Отсюда:
\begin{enumerate}[label=\textup{(\alph*)}]
\item \emph{$F[\alpha]$~--- поле разложения $X^n-a$ тогда и только тогда, когда $X^n-1$ раскладывается на линейные
множители в $F[\alpha]$.} Если $X^n-1$ раскладывается в $F[\alpha]$, то все $\zeta_i\alpha$ лежат в $F[\alpha]$. Обратно,
если все корни $\zeta_i\alpha$ лежат в $F[\alpha]$, то и $\zeta_i=\zeta_i\alpha/\alpha\in F[\alpha]$ (здесь важно $\alpha\ne0$).
\item В частности, \emph{если $F$ содержит все корни $n$-й степени из $1$} (т.\,е.\ $X^n-1$ раскладывается в $F[X]$), то
$F[\alpha]$~--- поле разложения $X^n-a$. Обратное, вообще говоря, неверно (\rf{2.10a}).
\item Если $p=\operatorname{char}F\ne0$ и $n=p^kn'$ с $p\nmid n'$, то $X^n-1=(X^{n'}-1)^{p^k}$ (Фробениус), так что корни
$n$-й степени из $1$~--- это корни степени $n'$. В частности, $X^p-1=(X-1)^p$, и поле $F$ автоматически содержит все
корни $p$-й степени из $1$: при $a\notin F^p$ поле корня $F[\alpha]$ многочлена $X^p-a$~--- поле разложения (\rf{2.3}).
\end{enumerate}
Пусть $\operatorname{char}F\nmid n$. Тогда $X^n-1$ не имеет кратных корней: если $X^n-1=(X-\zeta)^2h$, то, дифференцируя,
получаем $nX^{n-1}=(X-\zeta)\bigl(2h+(X-\zeta)h'\bigr)$, откуда $n\zeta^{n-1}=0$ и $\zeta=0$, что невозможно. Значит, корни
$n$-й степени из $1$ в $E$ образуют подгруппу $\mu_n$ порядка $n$ в $E^\times$, она циклическая (упражнение 1-3, \rf{5.7c}),
и если $\zeta$~--- её образующая (\emph{первообразный} корень), то корни $X^n-a$~--- это $\alpha,\zeta\alpha,\dots,\zeta^{n-1}\alpha$, а
\emph{поле разложения $X^n-a$ равно $F[\alpha,\zeta]$}, и условие (a) означает $\zeta\in F[\alpha]$. (Если $a=0$, то
$X^n$ раскладывается уже над $F$, и $F[\alpha]=F$~--- поле разложения независимо от корней из $1$.)
\end{example}

\begin{example}[2.10a]
Поле $F[\alpha]$ может быть полем разложения $X^n-a$, хотя $F$ не содержит всех корней $n$-й степени из $1$~--- даже если
$X^n-a$ неприводим.

(1) $F=\Q$, $f=X^4+1$ ($n=4$, $a=-1$). Многочлен неприводим: $f(X+1)=X^4+4X^3+6X^2+4X+2$~--- многочлен Эйзенштейна при $p=2$
(\rf{1.16a}(b)).
Пусть $\alpha=e^{\pi i/4}$; корни $f$~--- это $\alpha,\alpha^3,\alpha^5,\alpha^7$, и все они лежат в $\Q[\alpha]$. Значит,
$\Q[\alpha]$~--- поле разложения, но $i=\alpha^2\notin\Q$. Здесь $\alpha$ сам является корнем из $1$ степени $8$, и
$\zeta=\alpha^2$~--- первообразный корень степени $4$ в $\Q[\alpha]$, как и требует \rf{2.10}(a) (см.\ рисунок ниже).
По модулю любого простого $p$ этот неприводимый над $\Q$ многочлен приводим, так что обращение \rf{1.18} неверно. При
$p=2$ имеем $X^4+1=(X+1)^4$. При нечётном $p$ корни $X^4+1$ лежат в поле разложения $X^8-1$ над $\F_p$, степень
которого равна порядку $p$ по модулю $8$ (\rf{2.10c}), т.\,е.\ не больше $2$ (так как $p^2\equiv1\pmod8$); поэтому каждый
неприводимый множитель $X^4+1$ над $\F_p$ имеет степень $\le2$ (ср.\ \rf{5.11j} для $X^4+1=\Phi_8$).

(2) $F=\F_7$, $f=X^9-2$. Пусть $g$~--- унитарный неприводимый делитель $f$ в $\F_7[X]$, $\alpha$~--- корень $g$ в поле
корня $\F_7[\alpha]=\F_7[X]/(g)$, и $q=7^k$~--- число элементов поля $\F_7[\alpha]$, $k=\deg g\le9$. В $\F_7^\times$ элемент $2$ имеет порядок $3$ ($2^3=8=1$),
поэтому $\alpha^{27}=2^3=1$, а $\alpha^9=2\ne1$; значит, порядок $\alpha$ в $\F_7[\alpha]^\times$ равен $27$. По теореме Лагранжа
$27\mid q-1=7^k-1$. Степени $7$ по модулю $27$ равны $7,22,19,25,13,10,16,4,1$, т.\,е.\ порядок $7$ по модулю $27$ равен
$9$; поэтому $9\mid k$, $k=9$, и $f$ неприводим. Далее, $\zeta=\alpha^3$ имеет порядок $9$, т.\,е.\ это первообразный корень
степени $9$ из $1$, лежащий в $\F_7[\alpha]$; по \rf{2.10}(a) $\F_7[\alpha]$~--- поле разложения $f$. Но $\F_7$ не содержит
первообразного корня степени $9$: $9\nmid|\F_7^\times|=6$.

(3) Приводимый пример: $F=\Q$, $f=X^3-1$, $\alpha=\omega$; $\Q[\omega]$~--- поле разложения, а $\omega\notin\Q$.

Поэтому попытка доказать <<$\zeta\in F$>>, записав $\zeta=\sum c_i\alpha^i$ и возведя в $n$-ю степень, обречена: в примере (1)
$\zeta=\alpha^2$, $\zeta^4=\alpha^8=1$, и получающийся многочлен $X^8-1$ действительно делится на минимальный многочлен
$X^4+1$ элемента $\alpha$~--- противоречия нет.
\end{example}

\par\noindent\begin{minipage}{\linewidth}
\begin{center}
\begin{tikzpicture}[scale=1.5]
\draw[->,gray] (-1.6,0)--(1.6,0); \draw[->,gray] (0,-1.4)--(0,1.4);
\draw (0,0) circle (1);
\foreach \k in {1,3,5,7} {\fill[cold] ({45*\k}:1) circle (1.6pt);}
\foreach \k in {0,2,4,6} {\fill[cnew] ({45*\k}:1) circle (1.6pt);}
\node[cold,above right] at (45:1) {$\alpha$};
\node[cold,above left] at (135:1) {$\alpha^3$};
\node[cold,below left] at (225:1) {$\alpha^5$};
\node[cold,below right] at (315:1) {$\alpha^7$};
\node[cnew,above right] at (90:1) {$i=\alpha^2$};
\node[cnew,below right] at (0:1) {$1=\alpha^8$};
\node[cnew,below left] at (180:1) {$-1=\alpha^4$};
\node[cnew,below right] at (270:1) {$-i=\alpha^6$};
\end{tikzpicture}
\end{center}
\figcap{Пример \rf{2.10a}(1): корни $X^4+1$~--- нечётные степени $\alpha=e^{\pi i/4}$ (синие); чётные степени $\alpha$~---
корни $X^4-1$ (красные). Все они лежат в $\Q[\alpha]$, в том числе $i=\alpha^2\notin\Q$.}
\end{minipage}\par\medskip

\begin{proposition}[2.10b]
Пусть $X^n-a\in F[X]$ неприводим, $\alpha$~--- его корень, и $F[\alpha]$~--- поле разложения $X^n-a$. Тогда $X^n-1$
раскладывается над $F$ в каждом из следующих случаев:
\begin{enumerate}[label=\textup{(\alph*)}]
\item $n=\ell$~--- простое число;
\item $\operatorname{char}F\nmid n$ и $\gcd\bigl(n,\varphi(n)\bigr)=1$, где $\varphi$~--- функция Эйлера (например, $n=15$,
$\varphi(15)=8$).
\end{enumerate}
Кроме того, если $F\subset\R$ и $\alpha$ вещественен, то $F[\alpha]$~--- поле разложения $X^n-a$ лишь при $n\le2$.
\end{proposition}

\begin{proof}
По \rf{2.10}(a) $X^n-1$ раскладывается в $F[\alpha]$, а $[F[\alpha]:F]=n$ по неприводимости.

(a) Если $\operatorname{char}F=\ell$, то $X^\ell-1=(X-1)^\ell$ раскладывается над $F$ всегда. Пусть
$\operatorname{char}F\ne\ell$, и $\zeta\ne1$~--- корень $X^\ell-1$ в $F[\alpha]$; он имеет порядок $\ell$, и все корни
$X^\ell-1$~--- его степени. Из $F\subset F[\zeta]\subset F[\alpha]$ по \rf{1.20} степень $[F[\zeta]:F]$ делит $\ell$. С другой
стороны, $\zeta$~--- корень многочлена $X^{\ell-1}+\dots+X+1$ степени $\ell-1$, так что $[F[\zeta]:F]\le\ell-1<\ell$.
Значит, $[F[\zeta]:F]=1$, $\zeta\in F$, и $X^\ell-1=\prod_{i}(X-\zeta^i)$ раскладывается над $F$.

(b) Пусть $\zeta\in F[\alpha]$~--- первообразный корень степени $n$ из $1$ (\rf{2.10}). Как и в (a), $[F[\zeta]:F]$ делит $n$.
Кроме того, $[F[\zeta]:F]$ делит $\varphi(n)$: группа Галуа расширения $F[\zeta]/F$ вкладывается в $(\Z/n\Z)^\times$ (\rf{5.8}(c); здесь используется теория Галуа
глав~3 и~5),
а её порядок равен степени. Так как $\gcd(n,\varphi(n))=1$, степень равна $1$.

Вещественный случай: $F[\alpha]\subset\R$, а вещественные корни из $1$~--- это только $\pm1$; если $n\ge3$, то $X^n-1$ имеет
невещественный корень и не раскладывается в $F[\alpha]$.
\end{proof}

В примерах \rf{2.10a} условие (b) нарушено: $\gcd(4,\varphi(4))=2$, $\gcd(9,\varphi(9))=\gcd(9,6)=3$. Типичный
пример к вещественному случаю: $\Q[\sqrt[3]2]$ не является полем разложения $X^3-2$.

С точки зрения теории Галуа поле разложения $F[\zeta,\alpha]$ (при $\operatorname{char}F\nmid n$) строится в два шага:
\emph{круговой} шаг $F\subset F[\zeta]$ (его степень делит $\varphi(n)$, \rf{5.8}(c)) и \emph{шаг Куммера}
$F[\zeta]\subset F[\zeta,\alpha]$ (его степень~--- наименьшее $m\ge1$ с $\alpha^m\in F[\zeta]$, и $m\mid n$; глава~5, \rf{5.27}).
Для неприводимого $X^n-a$ (как в \rf{2.10b}) равенство $F[\alpha]=F[\zeta,\alpha]$ означает, что круговой шаг <<спрятан>>
внутри присоединения $\alpha$, а это при $\gcd(n,\varphi(n))=1$ возможно лишь при тривиальном круговом шаге, т.\,е.\ при
$\zeta\in F$. Без неприводимости это неверно: для $X^5-32=(X-2)(X^4+2X^3+4X^2+8X+16)$ над $\Q$ и корня $\alpha=2\zeta_5$
($\zeta_5=e^{2\pi i/5}$) поле $\Q[\alpha]=\Q[\zeta_5]$~--- поле разложения, $\gcd(5,4)=1$, но $\zeta_5\notin\Q$; так что и в
\rf{2.10b}(a) условие неприводимости нужно.

\par\noindent\begin{minipage}{\linewidth}
\begin{center}
\begin{tikzpicture}[x=1cm,y=1cm]
\node (F) at (0,0) {$F$};
\node[cnew] (Z) at (-2.4,1.6) {$F[\zeta]$};
\node[cold] (A) at (2.4,1.6) {$F[\alpha]$};
\node[cnew] (ZA) at (0,3.2) {$F[\zeta,\alpha]$};
\draw (F) -- node[below left,font=\footnotesize]{делит $\varphi(n)$} (Z);
\draw (F) -- node[below right,font=\footnotesize]{$n$} (A);
\draw (Z) -- node[above left,font=\footnotesize]{делит $n$} (ZA);
\draw[cnew,thick] (A) -- node[above right,font=\footnotesize]{$1\iff\zeta\in F[\alpha]$} (ZA);
\end{tikzpicture}
\end{center}
\figcap{Поле разложения неприводимого $X^n-a$ при $\operatorname{char}F\nmid n$: круговой шаг $F\subset F[\zeta]$ и шаг Куммера
$F[\zeta]\subset F[\zeta,\alpha]$. Поле $F[\alpha]$~--- поле разложения ровно тогда, когда правый верхний шаг тривиален; тогда
$F[\zeta]\subset F[\alpha]$, степень кругового шага делит и $n$, и $\varphi(n)$, и при $\gcd(n,\varphi(n))=1$ он тривиален.}
\end{minipage}\par\medskip

\subsection*{Поле разложения $X^m-1$ над $\F_p$}

\begin{proposition}[2.10c]
Пусть $p$ простое, $m=p^kn$ с $p\nmid n$, и $d$~--- порядок $p$ по модулю $n$ (наименьшее $d\ge1$ с $p^d\equiv1\pmod n$).
Тогда поле разложения $E$ многочлена $X^m-1$ над $\F_p$ имеет степень $[E:\F_p]=d$ и состоит из $p^d$ элементов.
\end{proposition}

\begin{proof}
По \rf{2.10}(c) (для показателя $m=p^kn$) $X^m-1=(X^n-1)^{p^k}$, поэтому $X^m-1$ и $X^n-1$ имеют одни и те же поля разложения (\rf{2.4e}(1)), и можно
считать $m=n$. Так как $p\nmid n$, корни $X^n-1$ в $E$ образуют циклическую группу $\mu_n$ порядка $n$ (\rf{2.10}); пусть
$\zeta$~--- её образующая. Все корни~--- степени $\zeta$, поэтому $E=\F_p[\zeta]$. Пусть $e=[E:\F_p]$; тогда $|E|=p^e$.

$d\le e$: группа $E^\times$ имеет порядок $p^e-1$ и содержит элемент $\zeta$ порядка $n$, поэтому $n\mid p^e-1$ (теорема
Лагранжа), т.\,е.\ $p^e\equiv1\pmod n$, и по минимальности $d\le e$ (более того, $d\mid e$).

$e\le d$: рассмотрим $K=\{x\in E\mid x^{p^d}=x\}$. Так как $x\mapsto x^{p^d}$~--- гомоморфизм колец (степень эндоморфизма
Фробениуса), $K$ замкнуто относительно сложения и умножения, а вместе с $x\ne0$ содержит и $x^{-1}$; значит, $K$~---
подполе. Оно содержит $\F_p$ (по малой теореме Ферма $x^p=x$ при $x\in\F_p$) и $\zeta$: так как $n\mid p^d-1$, имеем
$\zeta^{p^d}=\zeta\cdot\zeta^{p^d-1}=\zeta$. Поэтому $K\supset\F_p[\zeta]=E$, т.\,е.\ все $p^e$ элементов $E$~--- корни многочлена
$X^{p^d}-X$ степени $p^d$. Следовательно, $p^e\le p^d$ и $e\le d$.
\end{proof}

Например, для $X^7-1$ над $\F_2$: $2^3=8\equiv1\pmod 7$, $d=3$, поле разложения состоит из $8$ элементов; соответственно
$X^7-1=(X+1)(X^3+X+1)(X^3+X^2+1)$ над $\F_2$ (кубические множители неприводимы, так как не имеют корней в $\F_2$). Для
$X^{12}-1$ над $\F_3$: $12=3\cdot4$, $X^{12}-1=(X^4-1)^3=(X-1)^3(X+1)^3(X^2+1)^3$, $d=\operatorname{ord}_4(3)=2$, поле
разложения из $9$ элементов. Для $X^{p^k}-1=(X-1)^{p^k}$ ($n=1$) поле разложения~--- само $\F_p$. Для $X^{p-1}-1$ имеем
$p\equiv1\pmod{p-1}$, $d=1$: многочлен раскладывается уже над $\F_p$, в согласии с малой теоремой Ферма (его корни~---
все $p-1$ элементов $\F_p^\times$). Предложение \rf{2.10c} с точки зрения теории Галуа~--- это \rf{5.8c}, а разложение кругового многочлена по модулю $p$~--- в \rf{5.11j}.

\begin{aside}[2.11]
Пусть $F$~--- поле. Для данного $n$ многочлены степени $n$ из $F[X]$, поле разложения которых имеет степень ровно $n!$,
могут существовать, а могут и не существовать~--- это зависит от $F$. Их нет при $n>1$, если $F=\CC$ (поле $\CC$
алгебраически замкнуто, \rf{5.6}, и поле разложения любого многочлена~--- само $\CC$), и при $n>2$, если $F=\R$ (поле
разложения изоморфно подполю $\CC$ по \rf{2.12}, так как $\CC$ расщепляет $f$ (\rf{5.6}), и потому имеет степень
$\le2<n!$) или $F=\F_p$ (группа Галуа конечного расширения $\F_p$
циклическая, \rf{4.20}, ср.\ \rf{4.25a}; она вкладывается в $S_n$, а при $n\ge3$ группа $S_n$ не циклическая, так что
её циклические подгруппы имеют порядок меньше $n!$). Над $\Q$ при $n=3$ такой
многочлен~--- $X^3+10X+1$ (\rf{2.6}); позже (\rf{4.32}) мы увидим, как выписать бесконечно много многочленов степени $n$
из $\Q[X]$ с полем разложения степени $n!$.
\end{aside}
